\documentclass[11pt]{article}

\usepackage[T1]{fontenc}
\usepackage[utf8]{inputenc}
\usepackage{lmodern}
\usepackage[margin=0.95in]{geometry}
\usepackage{amsmath,amssymb,amsthm,mathtools}
\usepackage{booktabs}
\usepackage{enumitem}
\usepackage{graphicx}
\usepackage{xcolor}
\usepackage{hyperref}
\usepackage{microtype}

\hypersetup{colorlinks=true, linkcolor=blue, citecolor=blue, urlcolor=blue}
\setlength{\parindent}{0pt}
\setlength{\parskip}{0.48em}
\setlist[itemize]{leftmargin=1.35em,itemsep=0.1em,topsep=0.15em}
\setlist[enumerate]{leftmargin=1.5em,itemsep=0.1em,topsep=0.15em}

\newtheorem{theorem}{Theorem}
\newtheorem{proposition}[theorem]{Proposition}
\newtheorem{corollary}[theorem]{Corollary}
\newtheorem{assumption}[theorem]{Assumption}
\newtheorem{definition}[theorem]{Definition}

\DeclareMathOperator{\Acc}{Acc}
\DeclareMathOperator{\AUPRC}{AUPRC}
\DeclareMathOperator{\AUROC}{AUROC}
\DeclareMathOperator{\AUV}{AUV}
\newcommand{\E}{\mathbb E}
\newcommand{\R}{\mathbb R}
\newcommand{\cc}{\mathrm{cc}}
\newcommand{\eps}{\varepsilon}
\newcommand{\method}{\textsc{LinkRadius}}
\newcommand{\Vuln}{\mathcal V}
\newcommand{\sys}{\mathcal S}
\newcommand{\threat}{\mathcal T}
\newcommand{\Load}{\mathsf{Load}}
\newcommand{\NotRun}{\textcolor{gray}{not run}}

\title{When Silent Agents Fail:\\
From Handoff Radii to System-Level Vulnerability in Latent Multi-Agent Systems}
\author{Draft Manuscript}
\date{\today}

\begin{document}
\maketitle

\begin{abstract}
Latent multi-agent systems allow language-model agents to exchange continuous
representations without exposing a natural-language transcript. This can make
coordination efficient, but it leaves designers without a calibrated answer to
two questions: which concrete handoff is closest to changing the final
decision, and how should those local weaknesses be summarized for an entire
recursive or non-recursive architecture?

We introduce \method{}, a hierarchical, attack-independent diagnostic. It
first audits whether each handoff carries example-specific information using
paired, answer-label-matched mismatched, zero, and moment-matched replacements.
It then uses antithetic isotropic probes to estimate the scale-normalized
end-to-end susceptibility of every terminal pairwise margin. Dividing the
clean margin by this susceptibility gives an empirical handoff radius: the
relative perturbation budget at which a sufficiently aligned direction is
locally expected to reach a decision boundary. We aggregate the minimum radius
over a declared set of admissible links into a budget-indexed system
vulnerability profile, \(\Vuln_{\sys,\mathcal D,\threat}(\eps)\in[0,1]\). Its
value is the predicted fraction of clean-correct examples with an allowed
one-shot failure boundary within budget \(\eps\), rather than an arbitrary
sigmoid or clipped gain.

We prove that isotropic directional derivatives recover squared end-to-end
margin susceptibility in expectation, derive a curvature-aware sufficient
condition for handoff stability, and lift it to one-shot system stability on
the exact unrolled computation graph. Finite probes do not verify the required
uniform curvature bound, so experimental profiles are explicitly empirical,
not certified. In the completed validation pilot, 13 of 40 GPQA examples were
clean-correct under deterministic forced-option scoring and remained eligible
on clean recapture. Across these examples, answer-label-matched relay
replacement reduced accuracy by 38.5 percentage points at both the first
planner-to-critic and critic-to-solver handoffs, while the first
solver-to-planner effect was 7.7 points. All 1,872 requested antithetic probe
pairs passed the frozen cast-quality checks, although a six-unit exact-gradient
diagnostic did not achieve finite-difference agreement and is not used as a
certificate.

On an exploratory, fixed execution-order prefix of six validation examples
(18 example--handoff units), the empirical radius achieved a realized-norm
crossed-threshold Spearman correlation of 0.750 and interval-censored
concordance of 0.740 under per-example PGD. The corresponding values were
0.316 and 0.753 for clean margin alone, and 0.600 and 0.523 for susceptibility
alone. Thus the preliminary evidence supports useful boundary ordering but not
uniform superiority to clean margin. Held-out test evaluation, transfer
attacks, system profiling, and protection experiments remain unobserved and
are required for confirmatory claims.
\end{abstract}

\section{Introduction}

Language-model agents increasingly communicate through hidden states, latent
memories, or key--value cache handoffs instead of fully decoded text
\cite{du2025interlat,zou2025latentmas,yang2026recursivemas}. These channels can
reduce token overhead and preserve information that would otherwise be lost
during textualization. Their opacity, however, changes the robustness problem.
A corrupted text message can be inspected; a corrupted latent handoff may
remain invisible until the system emits an incorrect answer.

Recent work has established that latent channels can carry attack-associated
signals and that edge-level handoff attacks can degrade latent multi-agent
systems \cite{wang2026outofsight,brito2026latentagents}. This establishes
possibility, but it does not answer the design questions faced before a
particular attack has been constructed:
\begin{quote}
\emph{Which latent handoff is closest to failure, at what budget, and how much
of the system's clean-correct behavior is exposed under the same threat model?}
\end{quote}

Three issues make these questions harder than measuring hidden-state movement.
First, a receiver may not use the sender's example-specific message at all.
Recent causal audits show that replacing or zeroing a latent relay can produce
a large effect even when replacing it with another example's relay produces
almost none \cite{cheng2026when,zhang2026channels}. Such a link is sensitive
to message presence but is weak evidence of inter-agent information transfer.
Second, a large terminal representation change need not cross the relevant
decision boundary, while a small change can flip an already narrow margin.
Third, ``system vulnerability'' is not a budget-free scalar. It depends on the
allowed links, norm, site-selection policy, task distribution, and clean
denominator. Squashing an uncalibrated propagation gain into \([0,1]\) does not
resolve this ambiguity.

We propose a three-level diagnosis. Level~0 asks whether a latent handoff carries
example-specific information by comparing the clean paired relay with a
mismatched-example relay, a zero relay, and moment-matched random noise. Stage~1
estimates how much small,
attack-independent random interventions change the final decision margins and
combines that susceptibility with the clean margin into a per-example,
per-handoff empirical radius. Level~2 aggregates the radii under an explicit
threat model. For a worst-site one-shot adversary, the architecture radius is
the smallest allowed handoff radius and its empirical cumulative distribution
is a system-level vulnerability curve.

The core decomposition is
\begin{equation}
    \underbrace{m_{x,j}}_{\text{clean margin}}
    \quad\text{versus}\quad
    \eps\,
    \underbrace{\chi_{x,e,j}}_{\text{system susceptibility}}
    \underbrace{A_{x,e,j}(v)}_{\text{attack alignment}},
    \label{eq:intro-decomposition}
\end{equation}
where $e$ is a concrete latent handoff, $j$ is a competing answer, and $v$ is
an attack direction. Random probes estimate $\chi$ without using $v$. A held-out
attack succeeds locally only when it is sufficiently large and sufficiently
aligned to overcome the available margin. This distinction explains why a
system can be intrinsically fragile even when a particular structured attack
has low success, and why random-noise ASR is not itself a sufficient
susceptibility measure.

At the system level, let \(\mathcal E_{\threat}(\sys,x)\) be the handoffs that
the declared threat model permits. We define
\begin{equation}
    \widehat\rho_{x,\sys,\threat}^{\min}
    =\min_{e\in\mathcal E_{\threat}(\sys,x)}
    \widehat\rho_{x,e}^{(1)},
    \qquad
    \widehat\Vuln_{\sys,\mathcal D,\threat}^{\mathrm{worst}}(\eps)
    =
    \Pr_{x\sim\mathcal D_{\cc}}
    [\widehat\rho_{x,\sys,\threat}^{\min}\le\eps].
    \label{eq:intro-system-profile}
\end{equation}
A value of \(0.70\) at \(\eps=0.05\), for example, means that 70\% of the
declared clean-correct population contains at least one allowed link whose
predicted local boundary is within a 5\% relative perturbation. We also report
random-site profiles so architectures with more handoffs are not compared only
through their minimum.

Our intended contributions, together with the current scope of evidence, are:
\begin{enumerate}
    \item We introduce a causal-use gate that separates example-specific latent
    communication from generic dependence on the presence or distribution of a
    relay.
    \item We define an attack-independent, scale-normalized susceptibility and
    empirical first-order radius on the exact unrolled handoff graph of a
    recursive or non-recursive latent multi-agent system.
    \item We derive an isotropic random-probe estimator and a curvature-aware
    margin-stability theorem, while sharply separating empirical prediction
    from certified robustness.
    \item We lift handoff radii into interpretable, budget-indexed \([0,1]\)
    vulnerability profiles with worst-site, random-site, fixed-policy, and
    causally-useful variants. These profiles support fair comparisons across
    horizons and architectures without hiding the threat model.
    \item We test whether local radii predict failure thresholds before using
    them to compare datasets, collaboration patterns, and prompt regimes. The
    current validation-only pilot is encouraging but mixed; held-out prediction
    and actionability remain planned rather than established contributions.
\end{enumerate}

\section{Problem Setting}
\label{sec:problem-setting}

\subsection{Latent systems as unrolled handoff graphs}

Let \(\sys\) be a latent multi-agent system with recursion horizon \(R\). We
represent one deterministic execution as an unrolled directed acyclic graph
\begin{equation}
    \mathcal G_{\sys,R}(x)
    =\bigl(\mathcal V_{\sys,R}(x),\mathcal E_{\sys,R}(x)\bigr),
\end{equation}
where a chronological edge \(e\in\mathcal E_{\sys,R}(x)\) transmits a latent
object \(z_e(x)\). Unrolling makes this representation valid for both acyclic
collaboration and systems that reuse the same roles over multiple rounds. Two
messages between the same roles at different rounds are distinct edges.

For each edge, define the downstream intervention map
\begin{equation}
    F_{x,e}:z\longmapsto \ell_{x,e}(z)\in\R^C,
    \label{eq:intervention-map}
\end{equation}
which replaces only \(z_e(x)\) by \(z\), recomputes every descendant of that
edge, and returns terminal option scores. All non-descendants, exogenous inputs,
model weights, and prompts remain fixed. The clean execution satisfies
\(F_{x,e}(z_e)=\ell_x\) for every edge, although its derivative with respect to
\(z_e\) is edge-specific.

The theorem-bearing setting uses deterministic, differentiable option scoring
and a fixed discrete control-flow trace during the local intervention. This is
implemented with forced scoring of each answer option rather than differentiating
through an \(\arg\max\) token. Architectures whose intervention changes a tool
call or other discrete branch are evaluated as behavioral extensions and are
not assigned the differentiable guarantee without a separate smooth surrogate.

\begin{definition}[One-shot latent-handoff threat model]
\label{def:threat-model}
A threat model
\(\threat=(\mathcal E_{\threat},\{\mathcal U_e\},\{s_{x,e}\},
\|\cdot\|_2,\pi)\)
declares (i) the eligible chronological handoffs,
(ii) a \(q_e\)-dimensional admissible perturbation subspace
\(\mathcal U_e\) at each edge, (iii) a positive reference scale \(s_{x,e}\)
and relative Euclidean/Frobenius norm, and (iv) a site-selection policy \(\pi\).
By default \(s_{x,e}=\|z_e\|_2\) for a nonzero relay; a predeclared clean
reference scale is required for a zero relay. An intervention changes one
eligible handoff once and then allows normal downstream execution. Worst-site
selection, uniform random-site selection, and a fixed-site policy are reported
separately.
\end{definition}

This internal intervention is a diagnostic gray/white-box abstraction of a
compromised RecursiveLink, latent adapter, cache-sharing transform, or transient
communication fault. It is not presented as a black-box capability to edit an
arbitrary hidden state. The attacker cannot change the question, prompts, base
models, final scorer, or other handoffs.

\subsection{Running recursive serial instantiation}

We study a recursive planner--critic--solver system. At recursion round $r$,
the system executes the ordered latent handoffs
\begin{align}
    p_r &= P_r(x,s_{r-1}), \\
    c_r &= C_r(x,p_r), \\
    s_r &= S_r(x,c_r),
    \label{eq:serial-system}
\end{align}
where $x$ denotes fixed exogenous input and $p_r,c_r,s_r$ are the transmitted
latent tensors. The final decoder maps $s_R$ to option scores
\begin{equation}
    \ell(x)=D_x(s_R)\in\R^C.
\end{equation}
The solver-to-planner message $s_r$ affects $p_{r+1}$; it does not feed the
planner in the same round.

The set of valid one-shot handoff sites at horizon $R$ is
\begin{equation}
\mathcal E_R=
\{(P{\to}C,r),(C{\to}S,r):1\le r\le R\}
\cup
\{(S{\to}P,r):1\le r<R\}.
\label{eq:edge-set}
\end{equation}
Thus $S{\to}P$ is inactive when $R=1$. The complete execution is an unrolled
graph of chronological micro-steps, even though role names repeat across
rounds. Mixture,
distillation, and deliberation patterns are represented by their own unrolled
edge sets rather than forced into this serial schedule.

\begin{figure}[t]
\centering
\fbox{\parbox[c][1.42in][c]{0.92\linewidth}{\centering
\textbf{Figure 1 placeholder: actual unrolled computation.}\\[0.4em]
Draw $P_r\rightarrow C_r\rightarrow S_r\rightarrow P_{r+1}$ for two rounds.
Highlight one handoff $e$, the paired/mismatched replacement gate, the
$\pm h$ antithetic probes, and the terminal option-margin readout.}}
\caption{The proposed diagnostic acts on a concrete transmitted latent tensor
in the actual serial execution. Downstream recomputation follows the descendants
of that handoff; no synchronous role-state approximation is used.}
\label{fig:unrolled-overview}
\end{figure}

\subsection{One-shot relative perturbations}

Let $z_e\in\R^{d_e}$ denote the vectorized latent object transmitted at an
eligible handoff $e\in\mathcal E_{\threat}(\sys,x)$. We inject once within
\(\mathcal U_e\),
\begin{equation}
    \widetilde z_e=z_e+\delta_e,
    \qquad
    \delta_e\in\mathcal U_e,
    \qquad
    \|\delta_e\|_2\le\eps s_{x,e},
    \label{eq:relative-perturbation}
\end{equation}
and recompute its descendants while holding the question, prompts, model
weights, decoding rule, and all exogenous inputs fixed. For matrix-valued or
cache-valued handoffs, $\|\cdot\|_2$ in Equation~\eqref{eq:relative-perturbation}
denotes the Frobenius norm after vectorization unless a restricted perturbation
subspace is stated explicitly.

Low-precision inference can make the requested perturbation differ from the
inserted perturbation. We therefore log
\begin{equation}
    \eps_{\mathrm{act}}
    =
    \frac{\|\operatorname{cast}(z_e+\delta_e)-
    \operatorname{cast}(z_e)\|_2}{s_{x,e}^{\mathrm{cast}}}
\end{equation}
where \(s_{x,e}^{\mathrm{cast}}\) is the cast-domain realization of the
declared reference scale. We use the actual post-cast norm in all analyses.
Probes that collapse to zero
or become substantially asymmetric after casting are resampled or excluded by
a pre-specified rule.

\subsection{Pairwise decision margins}

For a multiple-choice example $x$ with correct option $y$, let
\(F_{x,e,j}(z)\) be the length-normalized score of option $j$ under the
intervention map in Equation~\eqref{eq:intervention-map}. For each competitor
$j\ne y$, define the pairwise margin
\begin{equation}
    m_{x,e,j}(z)=F_{x,e,y}(z)-F_{x,e,j}(z).
    \label{eq:pairwise-margin}
\end{equation}
At the clean handoff, write \(m_{x,j}=m_{x,e,j}(z_e)\); its value is shared
across edges while its derivatives are not. An example is clean-correct when
every \(m_{x,j}>0\). Using all pairwise margins avoids assuming that the clean
runner-up remains the strongest competitor under perturbation.

For a clean-correct example, the exact one-shot radius of edge \(e\) is
\begin{equation}
    \rho^\star_{x,e}
    =
    \inf_{\substack{\delta\in\mathcal U_e:\\
    \min_{j\ne y}m_{x,e,j}(z_e+\delta)\le 0}}
    \frac{\|\delta\|_2}{s_{x,e}},
    \label{eq:exact-edge-radius}
\end{equation}
with the infimum defined as \(+\infty\) if the feasible set is empty. Ties
count as reaching the boundary. This is the target quantity: the smallest
relative intervention at that exact handoff that reaches any terminal decision
boundary. Computing it exactly is generally intractable. \method{} estimates
its local first-order analogue without using the attack later employed for
evaluation.

Our theorem-bearing evaluation uses tasks with explicit option scores, with
GPQA as the first pilot and MedQA as the planned replication. Free-form tasks
such as Math500 are included only after a continuous generative margin or
behavioral threshold has been validated separately; they do not inherit the
multiple-choice theorem by analogy.

\section{Stage 0: Is the Handoff Causally Used?}
\label{sec:causal-gate}

A robustness score is difficult to interpret if the receiver does not use the
sender's example-specific content. Before estimating fragility, we audit each
handoff with four replacements:
\begin{itemize}
    \item \textbf{Paired:} the clean relay $z_e(x)$ generated for the evaluated
    example.
    \item \textbf{Label-matched mismatch:} a relay $z_e(x')$ from a different
    example with the same correct option label, assigned by a within-label
    derangement, matched in shape, and rescaled to the paired relay norm. An
    arbitrary mismatch is retained as a supplementary control.
    \item \textbf{Zero:} an all-zero tensor with the same shape.
    \item \textbf{Moment-matched noise:} random noise matched to the empirical
    location and scale of clean relays at the same site.
\end{itemize}

The accuracy-based example-specific utility of link $e$ is
\begin{equation}
    U_e^{\mathrm{acc}}
    =
    \Acc_e^{\mathrm{paired}}
    -
    \Acc_e^{\mathrm{mismatched}},
    \label{eq:utility-accuracy}
\end{equation}
and its continuous counterpart is
\begin{equation}
    U_e^{\mathrm{margin}}
    =
    \E_x\!\left[
    m_x^{\mathrm{paired}}-m_x^{\mathrm{mismatched}}
    \right],
    \qquad
    m_x=\min_{j\ne y}m_{x,j}.
    \label{eq:utility-margin}
\end{equation}
The paired--zero contrast measures dependence on message presence, while the
paired--mismatched contrast measures dependence on example identity. A large
zeroing effect with negligible mismatched effect is therefore not treated as
evidence that the link transfers useful example-specific information.

The causal-use gate is evaluated with paired bootstrap intervals and a
pre-specified equivalence margin. It changes interpretation rather than
silently deleting attacker-accessible edges: the all-edge exposure profile
includes every predeclared accessible handoff, while a separate useful-link
profile includes only links whose identity effect passes a gate learned on
calibration data. A low-utility, high-fragility edge is reported as unnecessary
attack surface, not useful communication failure. If the current full-question
routing shows no identity effect, that result is retained and a disclosed,
pre-registered receiver-necessary setting---where the sender has information
unavailable to the receiver---is added as a causal positive control.

\section{LinkRadius: Attack-Independent Link Diagnosis}
\label{sec:method}

\subsection{End-to-end margin susceptibility}

For every clean-correct example, handoff, and competitor, define
\begin{equation}
    g_{x,e,j}
    =P_e\nabla_{z_e}m_{x,e,j}(z_e),
\end{equation}
where \(P_e\) is the orthogonal projector onto the declared admissible
subspace \(\mathcal U_e\). The scale-normalized end-to-end margin
susceptibility is
\begin{equation}
    \chi_{x,e,j}
    =
    s_{x,e}\,\|g_{x,e,j}\|_2.
    \label{eq:susceptibility}
\end{equation}
This quantity measures the largest first-order change in pairwise margin per
unit of relative latent perturbation. The end-to-end gradient already contains
the complete ordered downstream computation:
\begin{equation}
    g_{x,e,j}
    =
    P_eJ_{e\rightsquigarrow\mathrm{out}}^\top
    \nabla_{\mathrm{out}}m_{x,e,j},
    \label{eq:unrolled-jacobian}
\end{equation}
where $J_{e\rightsquigarrow\mathrm{out}}$ is the Jacobian product along the
actual descendants of $e$. Direct end-to-end measurement retains contractions,
cancellations, and the receiver's ability to ignore the relay. Products of
local Jacobian norms are used only as a mechanistic upper-bound analysis in the
appendix.

\subsection{Exact distance in the affine model}

Linearizing each pairwise margin within \(\mathcal U_e\) gives
\begin{equation}
    \widetilde m_{x,e,j}(\delta)
    =m_{x,j}+\langle g_{x,e,j},\delta\rangle.
    \label{eq:affine-margin}
\end{equation}
Define the pairwise first-order radius
\begin{equation}
    \rho^{(1)}_{x,e,j}
    =
    \begin{cases}
    m_{x,j}/\chi_{x,e,j}, & \chi_{x,e,j}>0,\\
    +\infty, & \chi_{x,e,j}=0,
    \end{cases}
\end{equation}
and \(\rho^{(1)}_{x,e}=\min_{j\ne y}\rho^{(1)}_{x,e,j}\).

\begin{proposition}[Exact affine-model LinkRadius]
\label{prop:affine-radius}
For a strict clean-correct example, \(\rho^{(1)}_{x,e}\) is exactly the
minimum relative norm \(\|\delta\|_2/s_{x,e}\) that reaches any pairwise
decision boundary in the affine model in Equation~\eqref{eq:affine-margin}.
For a competitor with nonzero \(g_{x,e,j}\), the minimum-norm direction is
\begin{equation}
    \delta_{x,e,j}^\star
    =-\frac{m_{x,j}}{\|g_{x,e,j}\|_2^2}g_{x,e,j}.
\end{equation}
\end{proposition}

Thus \(\rho^{(1)}\) is an exact geometric distance for the local affine
model, not the exact nonlinear radius \(\rho^\star\) in
Equation~\eqref{eq:exact-edge-radius}.

\subsection{Two-sided isotropic random probes}

For $u_k$ sampled uniformly from the unit sphere in the admissible perturbation
subspace \(\mathcal U_e\), define the two-sided response
\begin{equation}
    \widehat d_{h,k,j}
    =
    \frac{
    m_{x,e,j}(z_e+hs_{x,e}u_k)
    -
    m_{x,e,j}(z_e-hs_{x,e}u_k)
    }{2h}.
    \label{eq:central-difference}
\end{equation}
The squared susceptibility estimator is
\begin{equation}
    \widehat\chi_{h,x,e,j}^{\,2}
    =
    \frac{q_e}{K}
    \sum_{k=1}^{K}\widehat d_{h,k,j}^{\,2}.
    \label{eq:chi-estimator}
\end{equation}
The probes are generated independently of attack trajectories, failure labels,
and held-out test outcomes. Where automatic differentiation is available, the
exact gradient norm is used only as an estimator-validation baseline.

\begin{proposition}[Isotropic recovery and finite-probe variance]
\label{prop:isotropic-recovery}
Let $u$ be uniform on the unit sphere in the $q_e$-dimensional admissible
subspace, and let
\begin{equation}
    d_j(u)
    =
    \left.\frac{\partial}{\partial t}
    m_{x,e,j}(z_e+ts_{x,e}u)\right|_{t=0}.
\end{equation}
Then
\begin{equation}
    \E_u[d_j(u)^2]
    =
    \frac{\chi_{x,e,j}^2}{q_e}.
\end{equation}
For independent directions,
\begin{equation}
    Q_j=\frac{q_e}{K}\sum_{k=1}^K d_j(u_k)^2
\end{equation}
is unbiased and consistent for \(\chi_{x,e,j}^2\), with
\begin{equation}
    \operatorname{Var}(Q_j)
    =\chi_{x,e,j}^4
    \frac{2(q_e-1)}{K(q_e+2)}.
    \label{eq:probe-variance}
\end{equation}
If the third derivative of the one-dimensional margin trace is bounded by
\(B\) on \([-h,h]\), replacing \(d_j(u_k)\) with the central difference in
Equation~\eqref{eq:central-difference} incurs error at most \(Bh^2/6\).
\end{proposition}

The squared estimator is unbiased for exact directional derivatives;
\(\sqrt{Q_j}\) and the resulting minimum radius are not unbiased. The
large-\(q_e\) relative standard deviation of \(Q_j\) is approximately
\(\sqrt{2/K}\). The current validation protocol uses $K=8$ with repeated
probe seeds for both estimation and exploratory prediction; $K\in\{16,32\}$
is reserved for a probe-count ablation rather than silently substituted after
seeing outcomes. Finite \(h\), low-precision casting, and taking minima over
competitors and edges add further error. The implementation checks that
the realized positive and negative offsets satisfy predeclared norm,
antipodality, and cosine tolerances; otherwise it resamples or fits the
directional slope using realized offsets.

\subsection{Empirical first-order link radius}

Finite probes give the empirical first-order radius
\begin{equation}
    \widehat\rho^{(1)}_{h,x,e}
    =
    \min_{j\ne y}
    \begin{cases}
    m_{x,j}/\widehat\chi_{h,x,e,j}, &
    \widehat\chi_{h,x,e,j}>0,\\
    +\infty, & \widehat\chi_{h,x,e,j}=0.
    \end{cases}
    \label{eq:empirical-radius}
\end{equation}
After \(h\) is fixed on validation data, we suppress it in
\(\widehat\rho^{(1)}_{x,e}\) when no ambiguity arises.
At relative budget $\eps$, the dimensionless link load is
\begin{equation}
    \Load_{x,e}(\eps)
    =
    \frac{\eps}{\widehat\rho^{(1)}_{h,x,e}}.
    \label{eq:load-ratio}
\end{equation}
This load is useful for calibration but is unbounded; it is not the paper's
\([0,1]\) system score. A small radius, or equivalently a large load, means
that a sufficiently aligned direction can locally exhaust the affine margin at
a smaller budget.

\subsection{Curvature-aware stability}

\begin{assumption}[Local curvature]
\label{ass:curvature}
Fix a maximum audited relative radius \(\bar\eps>0\). For a pairwise margin
$m_{x,e,j}$, suppose it is twice continuously differentiable throughout
\(\{z_e+\delta:\delta\in\mathcal U_e,
\|\delta\|_2\le\bar\eps s_{x,e}\}\) and
\begin{equation}
    \sup_z
    \|P_e\nabla^2 m_{x,e,j}(z)P_e\|_{\mathrm{op}}
    \le H_{x,e,j}.
\end{equation}
Define $\kappa_{x,e,j}=s_{x,e}^2H_{x,e,j}$.
\end{assumption}

\begin{theorem}[Curvature-aware pairwise margin stability]
\label{thm:margin-stability}
Under Assumption~\ref{ass:curvature}, every perturbation satisfying
$\delta_e\in\mathcal U_e$, $\|\delta_e\|_2\le\eps s_{x,e}$, and
\(0\le\eps\le\bar\eps\) obeys
\begin{equation}
    m_{x,e,j}(z_e+\delta_e)
    \ge
    m_{x,j}
    -\eps\chi_{x,e,j}
    -\frac{1}{2}\eps^2\kappa_{x,e,j}.
    \label{eq:curvature-bound}
\end{equation}
Therefore the predicted option remains $y$ whenever the right-hand side is
positive for every $j\ne y$.
\end{theorem}

Define the pairwise sufficient root
\begin{equation}
    r^{\mathrm{cert}}_{x,e,j}
    =
    \begin{cases}
    \dfrac{-\chi_{x,e,j}+
    \sqrt{\chi_{x,e,j}^2+2\kappa_{x,e,j}m_{x,j}}}
    {\kappa_{x,e,j}}, & \kappa_{x,e,j}>0,\\[0.8em]
    m_{x,j}/\chi_{x,e,j}, & \kappa_{x,e,j}=0,\ \chi_{x,e,j}>0,\\
    +\infty, & \kappa_{x,e,j}=\chi_{x,e,j}=0.
    \end{cases}
    \label{eq:cert-radius}
\end{equation}
Then
\(\rho^{\mathrm{cert}}_{x,e}
=\min\{\bar\eps,\min_{j\ne y}r^{\mathrm{cert}}_{x,e,j}\}\)
satisfies \(\rho^\star_{x,e}\ge\rho^{\mathrm{cert}}_{x,e}\): every
\(\eps<\rho^{\mathrm{cert}}_{x,e}\) preserves the clean decision.

\begin{center}
\fbox{\parbox{0.91\linewidth}{
\textbf{Claim boundary.}
Equation~\eqref{eq:cert-radius} is a certificate only if the gradient norm and
a uniform curvature upper bound are conservatively verified throughout the
ball. Our finite-probe experiments do neither. We therefore call
$\widehat\rho^{(1)}$ an \emph{empirical first-order radius}, not a certificate.
$\Load\ge1$ means ``the budget reaches the estimated affine boundary for a
sufficiently aligned direction,'' not ``every attack must succeed.'' Two probe
scales diagnose nonlinearity; they do not verify a Hessian bound.}}
\end{center}

\subsection{Attack alignment explains residual variation}

For a unit-norm attack direction $v\in\mathcal U_e$, define its signed harmful
alignment with competitor $j$ as
\begin{equation}
    A_{x,e,j}(v)
    =
    -\frac{\langle g_{x,e,j},v\rangle}{\|g_{x,e,j}\|_2}
    \in[-1,1].
    \label{eq:alignment}
\end{equation}
We set \(A_{x,e,j}(v)=0\) when \(g_{x,e,j}=0\). Locally,
\begin{equation}
    m_{x,e,j}(z_e+\eps s_{x,e}v)
    =
    m_{x,j}
    -\eps\chi_{x,e,j}A_{x,e,j}(v)
    +O(\eps^2).
    \label{eq:alignment-decomposition}
\end{equation}
In the affine model, the threshold along \(v\) is
\begin{equation}
    \rho^{(1)}_{x,e}(v)
    =
    \min_{\substack{j\ne y:\ A_{x,e,j}(v)>0}}
    \frac{m_{x,j}}{\chi_{x,e,j}A_{x,e,j}(v)}
    \ge \rho^{(1)}_{x,e},
    \label{eq:directional-radius}
\end{equation}
with the minimum of an empty set defined as \(+\infty\).
Thus $\chi$ measures system susceptibility, $A(v)$ measures attack quality,
and $m$ measures remaining decision margin. Alignment is a post-hoc explanatory
variable, not part of the attack-independent diagnostic. It should explain why
two attacks with the same norm can have different success rates and why a
small predicted radius need not make an arbitrary DiffMean direction succeed.

\subsection{Communication utility versus fragility}

We summarize each link with its causal utility $U_e$ and its vulnerable
fraction at a validation-fixed reference budget,
\begin{equation}
    F_e(\eps_0)
    =\Pr_{x\sim\mathcal D_{\cc}}
    [\widehat\rho_{h,x,e}^{(1)}\le\eps_0].
\end{equation}
Unlike a reciprocal-median score, this quantity remains bounded when an
estimated radius is zero or infinite. Combining utility and fragility produces
four qualitatively different design cases:
\begin{center}
\small
\begin{tabular}{lll}
\toprule
& \textbf{Low fragility} & \textbf{High fragility} \\
\midrule
\textbf{High utility} & Valuable, robust relay & Critical link to protect \\
\textbf{Low utility} & Largely irrelevant relay & Unnecessary attack surface \\
\bottomrule
\end{tabular}
\end{center}
This map makes a null result useful: a link that moves hidden states but fails
the identity test should not be described as a valuable communication channel,
even if it appears sensitive to zeroing.

\section{From Link Radii to System Vulnerability Profiles}
\label{sec:system-profile}

\subsection{Exact budget-conditioned vulnerability}

Let \(\mathcal E_{\threat}(\sys,x)\) be the attacker-accessible handoffs
declared before evaluation. Under an adaptive one-shot site policy, define the
exact system radius
\begin{equation}
    \rho^\star_{x,\sys,\threat,\mathrm{adv}}
    =
    \min_{e\in\mathcal E_{\threat}(\sys,x)}\rho^\star_{x,e},
    \label{eq:exact-system-radius}
\end{equation}
with the minimum of an empty set equal to \(+\infty\). For a stated
clean-correct population \(\mathcal D_{\cc}\), the exact vulnerability
profile is
\begin{equation}
    \Vuln^\star_{\sys,\mathcal D,\threat,\mathrm{adv}}(\eps)
    =
    \Pr_{x\sim\mathcal D_{\cc}}
    [\rho^\star_{x,\sys,\threat,\mathrm{adv}}\le\eps].
    \label{eq:exact-vulnerability-profile}
\end{equation}
This is a cumulative distribution function and therefore lies in \([0,1]\)
and is nondecreasing in \(\eps\). It answers a precise question: what fraction
of initially correct examples admits some allowed handoff and some admissible
one-shot direction that reaches a decision boundary within budget \(\eps\)?
It is conditional on the dataset, horizon, interface metric, accessible links,
and site-selection policy; it is not an architecture-intrinsic constant.

Adaptive worst-site selection legitimately counts a larger set of handoffs as
a larger attack surface. To separate this exposure from average per-link
fragility, for any predeclared site distribution \(\pi(e\mid\sys,x)\) we also
define
\begin{equation}
    \Vuln^\star_{\sys,\mathcal D,\threat,\pi}(\eps)
    =
    \E_{x\sim\mathcal D_{\cc}}
    \left[
    \sum_{e\in\mathcal E_{\threat}(\sys,x)}
    \pi(e\mid\sys,x)
    \mathbf 1\{\rho^\star_{x,e}\le\eps\}
    \right].
    \label{eq:policy-vulnerability-profile}
\end{equation}
Uniform \(\pi\) gives a random-site profile. A fixed-site profile uses a
single site selected on validation data and held fixed for every test example.

\subsection{Empirical LinkRadius profile}

For \(n\) held-out clean-correct examples, the plug-in LinkRadius profile is
\begin{equation}
    \widehat\Vuln^{\mathrm{LR}}_{\sys,n,\threat,\mathrm{adv}}(\eps)
    =
    \frac{1}{n}\sum_{i=1}^{n}
    \mathbf 1\!\left\{
    \min_{e\in\mathcal E_{\threat}(\sys,x_i)}
    \widehat\rho^{(1)}_{h,x_i,e}
    \le\eps
    \right\}.
    \label{eq:empirical-vulnerability-profile}
\end{equation}
It is a predicted, linearized vulnerability profile. It equals neither the
exact profile in Equation~\eqref{eq:exact-vulnerability-profile} nor observed
ASR by definition; held-out attacks test those empirical relationships.

The main all-edge profile includes every predeclared attacker-accessible link.
We additionally report a useful-link profile restricted to a set
\(\mathcal E_{\mathrm{use}}\) selected using causal identity effects on
calibration data only. The distinction separates total latent-interface
exposure from fragility of communication that demonstrably carries
example-specific information.

\begin{corollary}[One-shot system stability]
\label{cor:system-stability}
Suppose Assumption~\ref{ass:curvature} and the required conservative quantities
hold for every competitor at every eligible handoff. Define
\begin{equation}
    \rho^{\mathrm{cert}}_{x,\sys,\threat}
    =
    \min_{e\in\mathcal E_{\threat}(\sys,x)}
    \rho^{\mathrm{cert}}_{x,e}.
\end{equation}
Then every allowed one-shot intervention with relative norm
\(\eps<\rho^{\mathrm{cert}}_{x,\sys,\threat}\) preserves the terminal
decision. Consequently, if
\(C_{\sys,\mathcal D,\threat}(\eps)
=\Pr[\rho^{\mathrm{cert}}_{x,\sys,\threat}>\eps]\), then
\begin{equation}
    \Vuln^\star_{\sys,\mathcal D,\threat,\mathrm{adv}}(\eps)
    \le 1-C_{\sys,\mathcal D,\threat}(\eps).
\end{equation}
This inequality does not hold when empirical first-order radii replace the
conservative radii.
\end{corollary}

\subsection{Curve summaries and clean-accuracy accounting}

The full vulnerability curve is primary. We summarize it by
\begin{equation}
    \eps_q=\inf\{\eps:\Vuln(\eps)\ge q\},
\end{equation}
especially \(\eps_{50}\). If a single \([0,1]\) summary is needed, we
predeclare \(\eps_{\max}\) on validation data and use
\begin{align}
    \AUV_{[0,\eps_{\max}]}
    &=\frac{1}{\eps_{\max}}
    \int_0^{\eps_{\max}}\Vuln(\eps)\,d\eps \\
    &=\E_x\!\left[
    \left(1-\frac{\rho^{\min}_x}{\eps_{\max}}\right)_+
    \right].
    \label{eq:auv}
\end{align}
where \(\rho_x^{\min}\) is the exact or empirical system radius corresponding
to the plotted profile.
Unlike an arbitrary sigmoid, this scalar has an explicit budget range; it is
the expected vulnerable fraction under a uniform budget on that range.

Conditioning only on clean-correct examples can make an inaccurate system look
robust. We therefore also report predicted robust accuracy on the original
evaluation population:
\begin{equation}
    \widehat{\mathrm{RA}}_{\sys}(\eps)
    =\frac{1}{|\mathcal D|}\sum_{x\in\mathcal D}
    \mathbf 1\{\sys(x)=y_x\}
    \mathbf 1\{\widehat\rho^{\min}_{x,\sys}>\eps\}.
    \label{eq:predicted-robust-accuracy}
\end{equation}
Architecture, horizon, and prompt comparisons use paired raw IDs and the common
intersection-clean-correct set as the primary denominator. Each system's own
clean-correct curve and clean accuracy are reported secondarily. All compared
systems use the same option scorer, realized relative-budget semantics, and
attacker access. Worst-site and random-site profiles are both required because
the minimum in Equation~\eqref{eq:exact-system-radius} mechanically gives
architectures with more accessible edges more opportunities to fail.

\section{Experimental Design}
\label{sec:experimental-design}

\subsection{Research questions}

The paper is organized around seven questions:
\begin{enumerate}
    \item[\textbf{RQ0}] Do the evaluated latent handoffs carry
    example-specific information?
    \item[\textbf{RQ1}] Do finite random probes recover a stable local
    susceptibility across probe radii, seeds, and probe counts?
    \item[\textbf{RQ2}] Does the empirical radius predict per-example flips,
    empirical failure thresholds, and the most vulnerable handoff?
    \item[\textbf{RQ3}] Does prediction transfer to attacks constructed
    independently of the probes and test examples?
    \item[\textbf{RQ4}] Do aggregated radii predict system-level vulnerability
    curves across recursion horizons, datasets, and collaboration patterns?
    \item[\textbf{RQ5}] Do amplifying and corrective prompts change
    vulnerability through clean margin, latent susceptibility, or relay use?
    \item[\textbf{RQ6}] Does combining communication utility with fragility
    support a better link-protection decision than fixed or random selection?
\end{enumerate}

\subsection{Systems, tasks, and scope}

The completed pilot uses the sequential-light RecursiveMAS implementation with
planner, critic, and solver roles at $R=2$, batch size one, and latent length
32. Deterministic forced-option scoring defines correctness: an example is
eligible exactly when every finite gold-versus-competitor margin is positive.
Free-form generation remains a diagnostic and cannot exclude an otherwise
valid forced-margin example. The final paper records the exact checkpoint,
role-specific weights, RecursiveLink parameters, precision, and code revision
actually used rather than inheriting model details from the official
RecursiveMAS release. Expansion to other horizons is conditional on completing
the held-out $R=2$ boundary test.

GPQA is the theorem-bearing primary task because all answer competitors have
well-defined option scores. MedQA is the planned multiple-choice replication.
Math500 and MBPPplus are secondary behavioral transfer tasks only after a
candidate-set, verifier, or token-level continuous score is independently
validated. They do not inherit the multiple-choice theorem by analogy.

The pilot perturbs the first valid occurrence of each handoff type:
$P{\to}C@1$, $C{\to}S@1$, and $S{\to}P@1$. The implementation labels these
zero-indexed sites \texttt{p2c@0}, \texttt{c2s@0}, and \texttt{s2p@0}; the
manuscript uses one-indexed round labels. This keeps the first comparison small
while preserving different downstream paths. Expansion to every valid
occurrence in Equation~\eqref{eq:edge-set} is a later study and will always
report the injection round rather than collapsing repeated role pairs.

After local validation, we profile four collaboration patterns---sequential,
mixture, distillation, and deliberation---using their exact unrolled graphs
\cite{yang2026recursivemas}. Native released systems differ in roles,
backbones, training, and compute, so their comparison is explicitly a
\emph{whole-system profile}, not causal evidence that topology alone determines
robustness. A causal topology comparison additionally uses matched backbones,
training data, option scoring, and approximately matched inference compute.
Deliberation is theorem-bearing only when the local tool trace is fixed; results
that change discrete tool selection are labeled behavioral.

\subsection{Admissible perturbation interfaces}

Probes and evaluated attacks must occupy the same declared subspace. We write
\begin{equation}
    \delta_e=\eps s_{x,e}B_eu,
    \qquad B_e^\top B_e=I_{q_e},
    \qquad \mathcal U_e=\operatorname{span}(B_e).
    \label{eq:subspace-parameterization}
\end{equation}
For variable-length relays, \(B_e\) is a fixed channel-wise or otherwise
aligned parameterization that broadcasts over valid positions; it is fitted or
declared without test outcomes. The estimator dimension is \(q_e\), not the
raw tensor size. Universal, DiffMean, PCA, PGD, and random directions all use
these coordinates. This avoids comparing full-space susceptibility with an
attack constrained to a different broadcast or channel subspace.

\subsection{Architecture and prompt expansion}

The broad sweep is deliberately nested rather than fully factorial:
\begin{enumerate}
    \item sequential-light GPQA at \(R\in\{1,2,4\}\), followed by all valid
    sites and \(R=1,\ldots,5\);
    \item frozen-protocol replication on MedQA;
    \item compatible mixture and distillation systems as whole-system external
    validity, followed by deliberation under fixed control flow;
    \item amplifying versus corrective prompts on sequential-light GPQA and one
    frozen MedQA replication.
\end{enumerate}
The prompt study is a mechanism analysis. On paired examples it decomposes
\begin{equation}
    \log\rho^{(1)}_{x,e,j}
    =\log m_{x,j}-\log\chi_{x,e,j},
    \label{eq:prompt-decomposition}
\end{equation}
and separately measures relay identity utility and the binding competitor.
Prompts are frozen before evaluation. They are not treated as an independent
proof of the theorem.

\subsection{Data and split discipline}

Raw example identifiers are partitioned before clean-correct filtering:
\begin{itemize}
    \item $\mathcal D_{\mathrm{atk}}$ (40\%) constructs universal, DiffMean,
    and PCA directions.
    \item $\mathcal D_{\mathrm{val}}$ (20\%) selects probe radius, probe count,
    attack budgets, and attack hyperparameters.
    \item $\mathcal D_{\mathrm{test}}$ (40\%) is evaluated once after the
    protocol is frozen.
\end{itemize}
No test example, test failure label, or test probe response is used to construct
an attack. Primary analyses use a single frozen clean reference and deterministic
option scoring. The raw partitions are identical across architectures, horizons,
and prompts. Mismatched-relay donors come from a fixed disjoint donor pool and
are matched by site, round, horizon, answer label, and tensor-shape or sequence-
length bucket. Comparisons across systems use paired raw IDs and the
intersection clean-correct set; system-specific conditional curves are reported
only as secondary views.

\subsection{Probe protocol}

The completed validation pilot screened all 40 validation examples. Thirteen
were forced-margin clean-correct (32.5\%) and all 13 retained that status on a
fresh clean recapture. The probe grid uses $K=8$ antithetic directions, three
fixed direction seeds (101, 202, and 303), two candidate relative radii
($10^{-3}$ and $3\times10^{-3}$), and the three first-round handoffs. This
produced 234 eligible example--edge--scale--seed configurations and 1,872
antithetic pairs; every pair passed the frozen post-cast acceptance rule. The
frozen empirical analysis permits a minimum effective count of six of eight
directions, although the completed pilot retained all eight in every eligible
configuration.

We record requested and realized perturbation norms, requested--realized
cosine, positive--negative antipodality, each pairwise margin response, the
binding competitor, clean rerun disagreement, and invalid outputs for every
probe. A six-unit exact-autograd diagnostic was finite but achieved zero of six
finite-difference agreement checks, with maximum relative error 1.043. We
therefore retain that result as a negative estimator diagnostic and make no
exact-gradient or certification claim. Increasing $K$ and resolving the
autograd discrepancy are follow-up measurements rather than hidden conditions
for reporting the empirical $K=8$ result.

\subsection{Independent attack families}

We evaluate the diagnostic against attacks with different alignment properties.
The currently executed validation experiment contains the last two families
below; the first three remain planned transfer tests:
\begin{itemize}
    \item \textbf{Held-out universal margin attack:} one direction per handoff
    type learned at $R=2$ on $\mathcal D_{\mathrm{atk}}$, then frozen and
    transferred to both $R=2$ and $R=4$ test conditions.
    \item \textbf{Held-out DiffMean:} a failure-associated direction estimated
    only from attack-training trajectories, with target construction and sign
    fixed in advance. We use target-conditioned or pairwise-margin construction
    rather than averaging heterogeneous wrong-answer directions that can cancel.
    \item \textbf{Held-out PCA direction:} a principal failure-associated
    subspace learned on attack-training data, evaluated without test adaptation.
    \item \textbf{Per-example PGD:} a near-worst-case white-box ceiling used to
    test whether the radius locates a local boundary; it is reported separately
    from transfer attacks.
    \item \textbf{Random direction:} a null attack baseline distinct from the
    probe directions used by the diagnostic.
\end{itemize}
Attacks are norm-normalized after casting. The current grid is
$\{3\times10^{-4},10^{-3},3\times10^{-3},10^{-2},3\times10^{-2},10^{-1}\}$.
PGD uses 20 optimization steps and independently targets every wrong option;
the minimum resulting pairwise margin defines failure. The empirical threshold
is the first pairwise-margin exit. Crossings between tested budgets are
interval-censored; examples that never cross are right-censored above the
maximum budget. Because independently optimized PGD runs can re-enter the safe
region at a larger requested budget, first crossing and interval-censored
ordering are primary, and raw non-monotone dose responses remain visible.

The random-direction control exposed a BF16-scale floor: all 18 rows at
$3\times10^{-4}$ failed the probe-derived direction-fidelity thresholds, while
none collapsed and their realized-to-requested norm ratios remained near one.
All 90 rows at $10^{-3}$ and above passed. We therefore retain the smallest
random cell as a numerical diagnostic, use $10^{-3}$ and above for the
direction-qualified random sensitivity analysis, and do not remove individual
rows based on outcomes.

\subsection{Baselines and ablations}

The complete radius is compared with:
\begin{itemize}
    \item clean margin $m$ alone;
    \item susceptibility $\widehat\chi$ alone;
    \item latent norm $\|z_e\|_2$ alone;
    \item number of downstream micro-steps or nominal path length;
    \item terminal hidden-state displacement under random noise;
    \item worst-of-$K$ random directional response;
    \item empirical random-noise flip rate on validation-fixed probes;
    \item the previous role-gain $\widehat\Lambda$ score, recomputed only where
    its required quantities are available;
    \item a factorized product of local gains on the correctly unrolled graph;
    \item the exact-gradient radius as a white-box measurement oracle.
\end{itemize}
Estimator ablations vary $K\in\{4,8,16,32\}$, one-sided versus antithetic
probing, probe radius, scale normalization, and terminal margin versus hidden
distance. Where feasible, exact autograd gradient norms provide a measurement
reference rather than an attack.

At the system level, we compare the LinkRadius profile with edge count,
recursion depth, minimum and average clean margin, average susceptibility,
path length, and aggregated versions of the previous \(\widehat\Lambda\)
score. Scores without budget units receive a monotone calibration fitted on
validation data before curve-calibration comparisons; ranking metrics may use
their raw values. The decisive ablation remains whether \(m/\widehat\chi\)
beats both \(m\) and \(\widehat\chi\) separately.

\subsection{Metrics and statistical protocol}

Primary metrics are:
\begin{itemize}
    \item per-example flip $\AUROC$ and $\AUPRC$ using
    $\eps/\widehat\rho_{h,x,e}^{(1)}$;
    \item Spearman correlation and interval-censored concordance between the
    predicted radius and observed attack threshold;
    \item within-example site-ranking accuracy and Kendall rank correlation;
    \item calibration of failure rate against normalized budget
    $\eps/\widehat\rho$;
    \item system-level error in \(\eps_{50}\), integrated absolute error between
    predicted and observed vulnerability curves, architecture rank agreement
    by \(\AUV\), and predicted versus observed robust accuracy;
    \item paired cluster-bootstrap confidence intervals with raw example as the
    resampling unit and probe/attack seeds nested within examples.
\end{itemize}
We also report clean accuracy, clean-clean flip rate when decoding is stochastic,
invalid-output rate, actual perturbation norm, and probe cost. Causal-use tests
use Holm correction across links and a predeclared equivalence margin. Probe
uncertainty is propagated with repeated direction sets in addition to the
example-level bootstrap.

\begin{table}[t]
\centering
\small
\setlength{\tabcolsep}{3.6pt}
\caption{Planned experimental sequence. Scaling is conditional on the preceding
gate, so compute is spent only after the diagnostic is identifiable.}
\label{tab:experiment-sequence}
\begin{tabular}{p{0.16\linewidth}p{0.35\linewidth}p{0.37\linewidth}}
\toprule
\textbf{Stage} & \textbf{Design} & \textbf{Decision} \\
\midrule
0: Causal use &
Paired, mismatched, zero, and moment-matched relays on GPQA; $R\in\{2,4\}$ &
Label every link by identity utility; if all full-routing links fail, retain the
negative result and add the disclosed receiver-necessary positive control. \\
\midrule
1: Estimator &
Two probe radii; $K=8$ pilot; exact-gradient check where feasible &
Proceed if realized probes survive casting and susceptibility is stable across
radii/seeds. \\
\midrule
2: Prediction &
Frozen GPQA test set; all valid sites, rounds, budgets, and independent attacks &
Proceed if the full radius beats both margin-only and susceptibility-only with
paired intervals. \\
\midrule
3: Replication &
MedQA replication at two horizons with two transfer attacks &
Test cross-domain stability without reopening hyperparameters. \\
\midrule
4: Profiling &
Worst-site, random-site, and useful-link curves across horizons and compatible
architectures &
Proceed only if local radii predict held-out thresholds; compare both common
and system-specific clean sets. \\
\midrule
5: Mechanism &
Amplifying versus corrective prompts with margin--susceptibility decomposition &
Explain profile movement without treating prompts as theorem validation. \\
\midrule
6: Actionability &
Protect lowest-radius link versus random, fixed-last, and highest-radius links &
Test whether diagnosis improves robustness at matched intervention cost. \\
\bottomrule
\end{tabular}
\end{table}

\section{Results}
\label{sec:results}

Results in this section have two explicitly separated statuses. RQ0 and RQ1
use the complete 40-row validation pilot and its 13 forced-margin clean-correct
examples. RQ2 uses an exploratory, fixed execution-order prefix of six of those
13 examples (18 example--handoff units). The prefix contains every frozen
budget at all three handoffs for both PGD and the random control; it was chosen
to obtain an early matched analysis while the remaining validation attack jobs
were incomplete. Probe seeds are averaged equally and are not counted as
independent examples. No held-out test outcomes have been accessed. RQ3--RQ6
remain planned unless a row is explicitly marked as completed validation
evidence.

\subsection{RQ0: Causal use of the latent handoffs}

Table~\ref{tab:causal-use} reports the complete validation causal audit on the
13 stable clean-correct examples. The paired--mismatched contrast is the
decisive identity quantity; paired--zero is a message-presence diagnostic.
All 360 scheduled causal-grid tasks completed: 117 yielded eligible sample
rows (39 per handoff and 39 per intervention mode), while 243 were
predeclared frozen-cohort exclusions rather than runtime failures.

\begin{table}[t]
\centering
\scriptsize
\setlength{\tabcolsep}{2.2pt}
\caption{Completed validation causal-use audit ($n=13$). Effects are paired
baseline minus intervention; 95\% paired-bootstrap intervals are in
parentheses. Paired accuracy is one by construction on the stable
forced-margin cohort. Moment-matched noise and private-evidence routing have
not yet been run.}
\label{tab:causal-use}
\resizebox{\linewidth}{!}{%
\begin{tabular}{lrrrrrr}
\toprule
Handoff & $n$ & Paired Acc. & $\Delta\Acc_{\mathrm{id}}$ &
$\Delta m_{\mathrm{id}}$ & $\Delta\Acc_{0}$ & $\Delta m_{0}$ \\
\midrule
$P{\to}C@1$ & 13 & 1.000 & 0.385 (0.154, 0.615) &
0.029 ($-0.148$, 0.178) & 0.308 (0.077, 0.538) &
0.025 ($-0.123$, 0.143) \\
$C{\to}S@1$ & 13 & 1.000 & 0.385 (0.154, 0.615) &
0.043 ($-0.111$, 0.189) & 0.231 (0.000, 0.462) &
0.044 ($-0.100$, 0.161) \\
$S{\to}P@1$ & 13 & 1.000 & 0.077 (0.000, 0.231) &
0.026 ($-0.079$, 0.123) & 0.077 (0.000, 0.231) &
$-0.040$ ($-0.245$, 0.114) \\
\bottomrule
\end{tabular}}
\end{table}

Answer-label-matched replacement reduced paired accuracy by 38.5 percentage
points at both $P{\to}C@1$ and $C{\to}S@1$; their bootstrap intervals exclude
zero. The corresponding $S{\to}P@1$ effect was 7.7 points with an interval
touching zero. Continuous minimum-margin effects were small and all intervals
crossed zero. These results are preliminary evidence that the first two
handoffs carry example-specific information on this cohort, while the evidence
for the first solver-to-planner handoff is weak. They do not yet establish the
stronger private-evidence positive control.

\subsection{RQ1: Probe validity, efficiency, and local linearity}

The finite-probe pipeline was numerically usable, but the exact-gradient
comparison was not. Table~\ref{tab:probe-validation} reports both facts rather
than treating a successful forward pass as gradient validation.

\begin{table}[t]
\centering
\small
\setlength{\tabcolsep}{4.0pt}
\caption{Completed numerical and probe-quality checks on validation. A
configuration is one example--edge--scale--seed cell; each contains eight
antithetic pairs. Exact-gradient agreement is a negative diagnostic, so the
reported radius remains empirical rather than certified.}
\label{tab:probe-validation}
\begin{tabular}{p{0.27\linewidth}p{0.18\linewidth}p{0.45\linewidth}}
\toprule
Check & Units & Result \\
\midrule
Screening numerics & 40 raw examples & 40/40 finite scorer rows and
40/40 finite relay interfaces \\
Forced-margin eligibility & 40 raw examples & 13/40 clean-correct (32.5\%) \\
Fresh clean recapture & 13 eligible examples & 13/13 retained; zero demotions
and zero promotions \\
$K=8$ antithetic probes & 720 grid tasks (234 eligible) & All 720 tasks
completed; 486 were frozen-cohort exclusions; 1,872/1,872 eligible pairs were
accepted, with eight directions in every eligible configuration \\
Exact-autograd diagnostic & 6 example--edge units & 6/6 finite, but 0/6 met
the finite-difference agreement rule; maximum relative error 1.043; retained
as a diagnostic outside the frozen probe gate because its evidence rows lacked
resolved system-identity metadata \\
\bottomrule
\end{tabular}
\end{table}

The failed exact-gradient agreement prevents a claim that the finite probes
recover the exact gradient norm in this implementation. It does not make the
subsequent empirical prediction test circular: the probes remain independent
of the PGD outcomes. The present claim is therefore narrower---the resulting
finite-probe score has empirical predictive value if it orders independently
measured failures.

\subsection{RQ2: Predicting failure boundaries}

The core test asks whether a pre-attack score predicts empirical attack
thresholds. The currently completed matched prefix contains six validation
examples, three handoffs, three probe seeds, and the full six-budget PGD grid.
The observed threshold is the first budget at which a pairwise margin exits the
correct region, with interval-censoring between adjacent budgets and
right-censoring when no tested budget crosses it. This is an exploratory
validation analysis, not the confirmatory held-out result.

\begin{table}[t]
\centering
\small
\setlength{\tabcolsep}{3.2pt}
\caption{Exploratory PGD prediction on six validation examples (18
example--handoff units). Flip metrics are descriptive means over the six fixed
budgets; threshold and site metrics use actual post-cast norms. Probe seeds are
averaged equally. No confidence intervals are claimed at this sample size.}
\label{tab:core-prediction}
\begin{tabular}{lrrrrr}
\toprule
Score & Flip $\AUROC$ & Flip $\AUPRC$ & Threshold $\rho_s$ & C-index & Site rank acc. \\
\midrule
Clean margin only & 0.770 & 0.736 & 0.316 & 0.753 & 0.333 \\
Susceptibility only & 0.523 & 0.484 & 0.600 & 0.523 & 0.556 \\
\method{} radius & 0.745 & 0.689 & 0.750 & 0.740 & 0.444 \\
\bottomrule
\end{tabular}
\end{table}

\begin{table}[t]
\centering
\scriptsize
\setlength{\tabcolsep}{3.0pt}
\caption{Per-budget PGD flip prediction. Each cell averages the metric equally
over three probe seeds; each seed evaluates the same 18 example--handoff units.
Boldface only identifies the largest point estimate and does not denote
statistical significance.}
\label{tab:pgd-budget-prediction}
\resizebox{\linewidth}{!}{%
\begin{tabular}{rrrrrrr}
\toprule
$\eps$ & \multicolumn{3}{c}{$\AUROC$} & \multicolumn{3}{c}{$\AUPRC$} \\
\cmidrule(lr){2-4}\cmidrule(lr){5-7}
& \method{} & Margin & Suscept. & \method{} & Margin & Suscept. \\
\midrule
$3\!\times\!10^{-4}$ & \textbf{0.867} & 0.800 & 0.662 &
\textbf{0.769} & 0.767 & 0.516 \\
$10^{-3}$ & \textbf{0.810} & 0.753 & 0.563 &
\textbf{0.744} & 0.742 & 0.519 \\
$3\!\times\!10^{-3}$ & \textbf{0.792} & 0.753 & 0.576 &
0.732 & \textbf{0.742} & 0.526 \\
$10^{-2}$ & 0.713 & \textbf{0.763} & 0.454 &
0.713 & \textbf{0.764} & 0.510 \\
$3\!\times\!10^{-2}$ & 0.694 & \textbf{0.875} & 0.523 &
0.670 & \textbf{0.833} & 0.484 \\
$10^{-1}$ & 0.597 & \textbf{0.675} & 0.359 &
0.504 & \textbf{0.567} & 0.351 \\
\bottomrule
\end{tabular}}
\end{table}

\begin{table}[t]
\centering
\small
\setlength{\tabcolsep}{3.2pt}
\caption{Threshold-ordering sensitivity to the budget coordinate. Actual
post-cast norms are primary; requested-grid values are a complete-grid
sensitivity analysis. ``Site pair'' is within-example concordance over
order-comparable handoff pairs.}
\label{tab:boundary-coordinate}
\resizebox{\linewidth}{!}{%
\begin{tabular}{llrrrr}
\toprule
Coordinate & Score & Threshold $\rho_s$ & C-index & Site pair & Site top-1 \\
\midrule
Realized & Clean margin & 0.316 & 0.753 & 0.500 & 0.333 \\
Realized & Susceptibility & 0.600 & 0.523 & 0.889 & 0.556 \\
Realized & \method{} & \textbf{0.750} & 0.740 & \textbf{0.889} & 0.444 \\
Requested & Clean margin & 0.203 & 0.731 & 0.500 & 0.667 \\
Requested & Susceptibility & 0.569 & 0.538 & \textbf{0.867} & \textbf{0.889} \\
Requested & \method{} & \textbf{0.786} & \textbf{0.752} & \textbf{0.867} & \textbf{0.889} \\
\bottomrule
\end{tabular}}
\end{table}

\begin{table}[t]
\centering
\scriptsize
\setlength{\tabcolsep}{3.2pt}
\caption{Observed failures by requested budget in the exploratory matched
prefix. Entries are failures out of six examples. Random results at
$3\times10^{-4}$ are omitted from the qualified control because all 18
example--handoff rows failed direction-fidelity checks without collapsing.}
\label{tab:dose-response-counts}
\resizebox{\linewidth}{!}{%
\begin{tabular}{rrrrrrr}
\toprule
& \multicolumn{3}{c}{Per-example PGD} &
\multicolumn{3}{c}{Qualified random direction} \\
\cmidrule(lr){2-4}\cmidrule(lr){5-7}
$\eps$ & $P{\to}C@1$ & $C{\to}S@1$ & $S{\to}P@1$ &
$P{\to}C@1$ & $C{\to}S@1$ & $S{\to}P@1$ \\
\midrule
$3\!\times\!10^{-4}$ & 2/6 & 1/6 & 2/6 & n/a & n/a & n/a \\
$10^{-3}$ & 2/6 & 3/6 & 2/6 & 2/6 & 2/6 & 0/6 \\
$3\!\times\!10^{-3}$ & 3/6 & 2/6 & 2/6 & 0/6 & 1/6 & 1/6 \\
$10^{-2}$ & 2/6 & 3/6 & 3/6 & 2/6 & 1/6 & 1/6 \\
$3\!\times\!10^{-2}$ & 2/6 & 2/6 & 2/6 & 1/6 & 1/6 & 1/6 \\
$10^{-1}$ & 3/6 & 2/6 & 2/6 & 2/6 & 2/6 & 1/6 \\
\bottomrule
\end{tabular}}
\end{table}

Across the five direction-qualified random budgets, PGD produced 35 failures
in 90 repeated example--edge--budget cells (38.9\%), compared with 18/90
(20.0\%) for the random control. This is descriptive rather than an
independent-sample test, but it confirms that optimized attacks find more
failures than an unrelated direction. Raw PGD failure counts are non-monotone
because each budget is optimized independently; the first-crossing threshold
retains any adversarial point already found inside the nested budget set.
Across all six requested random budgets, 90/108 rows passed the frozen
post-cast direction-quality rule. The 18 failures were exactly the six examples
at all three handoffs at $3\times10^{-4}$; all had cosine below 0.627 and
off-direction relative error above 0.779, while none collapsed and realized
norms remained approximately equal to the requested norms.

The central result is mixed. \method{} gives much stronger crossed-threshold
ordering than clean margin (0.750 versus 0.316 on realized norms), while clean
margin is slightly better on interval-censored concordance (0.753 versus
0.740). \method{} has the largest flip AUROC at the three smallest budgets,
but margin is stronger at the three largest. Susceptibility alone is generally
weaker on flip prediction and censored concordance, although it ties or exceeds
\method{} on some low-support site-ranking metrics. The evidence therefore
supports the narrower statement that \method{} adds useful local ordering,
especially near small budgets; it does not yet support uniform dominance over
margin.

\begin{figure}[t]
\centering
\fbox{\parbox[c][1.45in][c]{0.92\linewidth}{\centering
\textbf{Figure 2 placeholder: preliminary boundary prediction.}\\[0.4em]
Left: predicted $\widehat\rho$ versus interval-censored realized PGD threshold
for the six-example validation prefix. Right: per-budget AUROC for \method{},
clean margin, and susceptibility, with validation-only status marked.}}
\caption{The preliminary result favors \method{} for crossed-threshold ordering
and low-budget flip prediction, but not for every baseline comparison. The
final figure must add raw-example bootstrap intervals and held-out test data.}
\label{fig:radius-calibration}
\end{figure}

\subsection{RQ3: Transfer across independent attacks}

Per-example PGD tests the local worst-case interpretation of the radius. The
universal, DiffMean, and PCA directions test the harder claim that susceptibility
rankings remain useful for attacks not used in diagnosis. We report these
separately because a universal direction can fail through poor alignment even
when the underlying link is fragile.

\begin{table}[t]
\centering
\small
\setlength{\tabcolsep}{3.4pt}
\caption{Current transfer status. The PGD row repeats the exploratory
validation summary above; independent transfer attacks have not been run.
Alignment will be measured only after each attack is frozen and will explain
residuals rather than construct the attack-independent radius.}
\label{tab:attack-transfer}
\begin{tabular}{llrrrr}
\toprule
Attack & Status & Flip $\AUROC$ & Threshold $\rho_s$ & Site top-1 & Alignment \\
\midrule
Per-example PGD & Validation prefix & 0.745 & 0.750 & 0.444 & \NotRun \\
Random direction & Qualified control only & \NotRun & \NotRun & \NotRun & \NotRun \\
Universal margin & Planned & \NotRun & \NotRun & \NotRun & \NotRun \\
Held-out DiffMean & Planned & \NotRun & \NotRun & \NotRun & \NotRun \\
Held-out PCA & Planned & \NotRun & \NotRun & \NotRun & \NotRun \\
\bottomrule
\end{tabular}
\end{table}

\paragraph{Interpretation rule.}
Strong PGD prediction with weaker universal-attack prediction supports the
radius as a susceptibility diagnostic but not a universal ASR predictor. If
adding measured alignment substantially explains the residual, the intended
margin--susceptibility--alignment decomposition is supported. The paper must
not relabel poor attack alignment as a failure of the system-level diagnostic.
At present there is no RQ3 transfer result: the random direction is only a null
control, and no universal, DiffMean, or PCA attack has been evaluated.

\subsection{RQ4: System profiles across sites, horizons, and architectures}

RQ4 has not been run. In particular, the current six-example validation prefix
must not be promoted to an architecture-level vulnerability profile. The
tables and figure in this subsection specify the planned analysis that becomes
eligible only after the full validation grid is frozen and the held-out RQ2
result is evaluated.

Site is indexed by both role pair and exact injection round. We therefore avoid
averaging all $P{\to}C$ events into one value when they have different numbers
of downstream micro-steps. $S{\to}P$ results omit invalid $R=1$ conditions.

\begin{table}[t]
\centering
\small
\setlength{\tabcolsep}{3.5pt}
\caption{Planned condition-level radius and empirical $\eps_{50}$ analysis
across the unrolled graph. No entries in this table have yet been evaluated.}
\label{tab:site-horizon}
\begin{tabular}{llrrrrr}
\toprule
Dataset & Site & Horizon & Median $\widehat\rho$ & PGD $\eps_{50}$ & Univ. $\eps_{50}$ & Rank match \\
\midrule
GPQA & $P{\to}C@1$ & 2 & \NotRun & \NotRun & \NotRun & \NotRun \\
GPQA & $C{\to}S@1$ & 2 & \NotRun & \NotRun & \NotRun & \NotRun \\
GPQA & $S{\to}P@1$ & 2 & \NotRun & \NotRun & \NotRun & \NotRun \\
GPQA & $P{\to}C@1$ & 4 & \NotRun & \NotRun & \NotRun & \NotRun \\
GPQA & $C{\to}S@1$ & 4 & \NotRun & \NotRun & \NotRun & \NotRun \\
GPQA & $S{\to}P@1$ & 4 & \NotRun & \NotRun & \NotRun & \NotRun \\
MedQA & All valid & 2,4 & \NotRun & \NotRun & \NotRun & \NotRun \\
\bottomrule
\end{tabular}
\end{table}

We do not assume that more rounds monotonically increase fragility. Extra
computation may amplify, contract, overwrite, or ignore a relay. The empirical
claim is instead that \method{} tracks these horizon-specific changes better
than nominal path length.

We next aggregate exact-site radii into the profiles in
Section~\ref{sec:system-profile}. Predicted profiles are compared with empirical
PGD curves as the closest observed local worst-case reference and with held-out
universal curves as the transfer test.

\begin{figure}[t]
\centering
\fbox{\parbox[c][1.55in][c]{0.92\linewidth}{\centering
\textbf{Figure 3 placeholder: system vulnerability profiles.}\\[0.4em]
Left: worst-site, validation-fixed-site, and uniform-random-site
\(\widehat\Vuln^{\mathrm{LR}}(\eps)\) for sequential horizons. Middle:
predicted versus observed PGD curves. Right: whole-system profiles for
sequential, mixture, distillation, and fixed-trace deliberation.}}
\caption{Budget-indexed profiles preserve the meaning of a \([0,1]\)
vulnerability score. Worst-site curves include the number of accessible links;
uniform-site curves expose average per-link fragility.}
\label{fig:system-vulnerability-curves}
\end{figure}

\begin{table}[t]
\centering
\scriptsize
\setlength{\tabcolsep}{2.7pt}
\caption{System-level vulnerability profiling. \(\eps_0\) and
\(\eps_{\max}\) are frozen on validation data and shared across compared
systems. Architecture rows use the intersection-clean-correct set; system-own
conditional results and predicted robust accuracy are reported alongside them.}
\label{tab:system-profiles}
\begin{tabular}{llrrrrrr}
\toprule
Dataset & System / horizon & Clean Acc. &
\(\Vuln_{\mathrm{adv}}(\eps_0)\) &
\(\Vuln_{\mathrm{unif}}(\eps_0)\) &
\(\eps_{50}\) & \(\AUV\) & \(\widehat{\mathrm{RA}}(\eps_0)\) \\
\midrule
GPQA & Sequential, $R=1$ & \NotRun & \NotRun & \NotRun & \NotRun & \NotRun & \NotRun \\
GPQA & Sequential, $R=2$ & \NotRun & \NotRun & \NotRun & \NotRun & \NotRun & \NotRun \\
GPQA & Sequential, $R=4$ & \NotRun & \NotRun & \NotRun & \NotRun & \NotRun & \NotRun \\
GPQA & Mixture & \NotRun & \NotRun & \NotRun & \NotRun & \NotRun & \NotRun \\
GPQA & Distillation & \NotRun & \NotRun & \NotRun & \NotRun & \NotRun & \NotRun \\
GPQA & Deliberation, fixed trace & \NotRun & \NotRun & \NotRun & \NotRun & \NotRun & \NotRun \\
MedQA & Frozen-protocol replication & \NotRun & \NotRun & \NotRun & \NotRun & \NotRun & \NotRun \\
\bottomrule
\end{tabular}
\end{table}

The architecture table is interpreted at two levels. Matched-backbone,
matched-compute variants support claims about collaboration structure. Native
released variants support only whole-system comparisons because role models,
training, and link interfaces differ. We test whether LinkRadius ordering
agrees with observed curve ordering and report integrated absolute curve error,
not merely visually separated curves.

\subsection{RQ5: Prompt-induced vulnerability mechanisms}

RQ5 has not been run. Prompt effects remain a conditional extension rather
than evidence supporting the current local-boundary result.

Amplifying and corrective prompts are evaluated on paired examples under the
same topology, weights, handoffs, and budgets. The analysis asks which factor
moves the radius rather than merely whether prompts change ASR.

\begin{table}[t]
\centering
\small
\setlength{\tabcolsep}{3.3pt}
\caption{Prompt-regime decomposition on the intersection-clean-correct set.
Report paired changes and confidence intervals. Radius decomposition is
competitor-specific before aggregation because the binding competitor may
change between prompts.}
\label{tab:prompt-decomposition}
\begin{tabular}{lrrrrrrr}
\toprule
Regime & Identity $U$ & Margin $m$ & Suscept. $\widehat\chi$ &
Median $\widehat\rho$ & $\Vuln_{\mathrm{adv}}(\eps_0)$ & $\eps_{50}$ & $\AUV$ \\
\midrule
Amplifying & \NotRun & \NotRun & \NotRun & \NotRun & \NotRun & \NotRun & \NotRun \\
Corrective & \NotRun & \NotRun & \NotRun & \NotRun & \NotRun & \NotRun & \NotRun \\
Paired difference & \NotRun & \NotRun & \NotRun & \NotRun & \NotRun & \NotRun & \NotRun \\
\bottomrule
\end{tabular}
\end{table}

A corrective prompt may improve the profile by widening the clean margin,
reducing susceptibility, or bypassing the relay. Only the first two represent
greater robustness while preserving the same communication function; the
causal-use audit distinguishes them from lost relay utility.

\subsection{RQ6: Utility--fragility map and targeted protection}

RQ6 has not been run. The protection experiment will be selected only from a
frozen held-out vulnerability analysis; the current validation prefix is not
used to choose a protected link.

\begin{figure}[t]
\centering
\fbox{\parbox[c][1.42in][c]{0.92\linewidth}{\centering
\textbf{Figure 4 placeholder: link utility--fragility map.}\\[0.4em]
Plot paired--mismatched margin utility on the horizontal axis and
\(F_e(\eps_0)=\Pr[\widehat\rho_{x,e}\le\eps_0]\) on the vertical axis.
Label each site, injection round, and horizon. Shade the ``critical link'' and
``unnecessary attack surface'' regions.}}
\caption{Utility and fragility answer different design questions. A link is a
protection priority only when it is both causally useful and easy to destabilize.}
\label{fig:utility-fragility}
\end{figure}

As a final actionability test, validation data select
\begin{equation}
    e^\star
    =\arg\max_{e:U_e\ge\delta_U}
    \Pr_{x\sim\mathcal D_{\mathrm{val}}}
    [\widehat\rho^{(1)}_{h,x,e}\le\eps_0],
    \label{eq:protected-link-selection}
\end{equation}
where \(\delta_U\) and \(\eps_0\) are predeclared. The cleanest
selection-only experiment applies an integrity shield to one link, making that
single handoff unavailable to the latent attacker while leaving every other
link accessible. An optional practical realization selectively textualizes or
redundantly recomputes the chosen relay. The same mechanism and cost are used
for every selection policy; this tests diagnosis, not defense novelty. The
attacker is then evaluated against the worst remaining eligible site.

\begin{table}[t]
\centering
\small
\setlength{\tabcolsep}{3.6pt}
\caption{Matched-cost link protection. Report clean utility and robustness so
selection cannot win by disabling communication.}
\label{tab:protection}
\begin{tabular}{lrrrrr}
\toprule
Protected-link policy & Clean Acc. & Utility retained & PGD $\eps_{50}$ &
Univ. $\eps_{50}$ & $\AUV$ \\
\midrule
No protection & \NotRun & \NotRun & \NotRun & \NotRun & \NotRun \\
Random valid link & \NotRun & \NotRun & \NotRun & \NotRun & \NotRun \\
Fixed last link & \NotRun & \NotRun & \NotRun & \NotRun & \NotRun \\
Largest-radius link & \NotRun & \NotRun & \NotRun & \NotRun & \NotRun \\
Smallest-radius useful link & \NotRun & \NotRun & \NotRun & \NotRun & \NotRun \\
\bottomrule
\end{tabular}
\end{table}

\section{Discussion}

\paragraph{What the current evidence establishes.}
The validation pilot establishes that the complete measurement path is
feasible on sequential-light GPQA at $R=2$: forced-option margins and relay
interfaces were finite, all 13 selected examples repeated their clean-correct
status, the three first-round handoffs admitted complete $K=8$ antithetic probe
sets, and PGD produced empirical boundary crossings. On the fixed six-example
attack prefix, \method{} substantially improved crossed-threshold ordering over
clean margin and susceptibility and improved low-budget flip AUROC over both.
The causally audited $P{\to}C@1$ and $C{\to}S@1$ links also showed sizeable
accuracy effects under answer-label-matched relay replacement. Together these
are meaningful preliminary evidence that an attack-independent local radius
contains information about later optimized failures.

\paragraph{What the current evidence does not establish.}
The same prefix does not show uniform dominance. Clean margin slightly exceeds
\method{} on realized interval-censored concordance and on the descriptive
mean of per-budget AUROC, while susceptibility ties or exceeds \method{} on
some site-ranking summaries with very small support. Exact-gradient agreement
also failed. Most importantly, all boundary results so far use validation
examples, only six independent raw examples enter the attack analysis, and no
test outcome has been opened. We therefore do not claim held-out prediction,
estimator correctness relative to autograd, transfer beyond PGD, or a validated
system-level profile.

\paragraph{What a confirmatory positive result would establish.}
If the frozen radius reproduces its boundary-ordering advantage on held-out
test examples with raw-example bootstrap intervals, then small random probes
reveal a system property useful before selecting an attack. Transfer to
universal and failure-associated directions would further show that the
fragility ordering is not specific to one optimization procedure. Agreement
between aggregated \method{} profiles and observed worst-site curves would
then justify testing the architecture-level \([0,1]\) score, while
diagnosis-guided shielding would test whether the profile supports an actual
design decision.

\paragraph{What a negative result would establish.}
If the causal-use gate fails, the evaluated full-routing links are not carrying
the example-specific information attributed to them; their sensitivity is
latent-interface exposure rather than communication fragility, and the
receiver-necessary condition becomes the causal positive control. If the full
radius fails while margin-only succeeds, output confidence dominates the
measured latent sensitivity. If PGD follows the radius but universal attacks do
not, alignment---rather than system susceptibility---is the limiting factor.
If local prediction fails, architecture-level curves are not promoted to the
headline merely because they separate systems. Each branch is more informative
than reporting small average excess ASR over many conditions.

\paragraph{Why the random-control floor is informative.}
The smallest random budget preserved its requested norm but not its requested
direction after low-precision casting. Because this occurred in all examples
and sites at exactly one predeclared budget, and because none of the
interventions collapsed, it is best treated as a measurement-resolution limit
rather than selective missingness. Excluding the entire $3\times10^{-4}$ cell
from the direction-qualified random comparison is outcome-independent and
leaves the PGD analysis unchanged.

\paragraph{What the system score does and does not mean.}
\(\Vuln_{\sys,\mathcal D,\threat}(\eps)\) is a budget-conditioned fraction,
not a universal security grade. Changing the allowed handoffs, perturbation
subspace, scale convention, site policy, horizon, or clean population changes
the question and therefore the score. The full curve is primary;
\(\eps_{50}\) and \(\AUV\) are compact summaries with explicit units or
budget ranges.

\paragraph{Why this is not another attack paper.}
Attacks are used as held-out instruments for testing a pre-attack diagnostic.
The contribution is the ability to localize and rank latent failure exposure,
not a new way to maximize damage. This makes the work complementary to latent
attack construction and latent-channel defenses.

\section{Related Work}
\label{sec:related-work}

\paragraph{Latent communication.}
InterLat and LatentMAS replace textual inter-agent messages with continuous
representations, while RecursiveMAS organizes latent communication into a
recursive multi-agent computation \cite{du2025interlat,zou2025latentmas,
yang2026recursivemas}. These systems motivate handoff-level analysis because
their communication object is continuous and their execution is composed over
multiple downstream modules.

\paragraph{Latent attacks and safeguards.}
Latent-only interventions can reactivate attack-associated behavior in latent
multi-agent systems, particularly at inter-agent handoffs
\cite{wang2026outofsight}, while tampered KV-cache collaboration exposes an
integrity risk invisible to text-only checks \cite{brito2026latentagents}.
LCGuard primarily studies reconstruction-based sensitive-information leakage
and learned transformations for safer KV sharing \cite{asif2026lcguard}.
Our work is complementary: we ask whether attack-independent measurements can
identify fragile links and budgets before choosing a defense or attack.

\paragraph{Causal audits of communication.}
Recent work distinguishes the effect of a message's presence from the effect
of its example-specific identity using matched, mismatched, zero, and randomized
relays \cite{cheng2026when,zhang2026channels}. We adopt this distinction as a
prerequisite for robustness analysis: perturbation sensitivity is meaningful as
communication fragility only when the link transmits task-relevant content.

\paragraph{Random directional estimation and robustness margins.}
Random-direction finite differences are a standard way to recover gradient
information without backpropagation \cite{nesterov2017random}. Attack-independent
margin and local-Lipschitz robustness scores also precede this work, notably
CLEVER \cite{weng2018clever}. We therefore do not claim the spherical moment
identity or Taylor margin bound alone as novel. The contribution is their
system-specific use at causally audited latent handoffs, exact unrolled
execution, held-out threshold prediction, budget-indexed architecture
aggregation, and diagnosis-guided link selection.

\section{Limitations}

\method{} is local, norm-specific, subspace-specific, and margin-specific. Its
first-order radius is exact only for the affine model. Finite random probes can
underestimate narrow sensitivity directions, and minimizing noisy estimates
over competitors and edges introduces selection bias. Repeated probes and
exact-gradient checks quantify this error but do not create a certificate. A
small empirical radius indicates opportunity for a sufficiently aligned
perturbation, not success of an arbitrary universal, DiffMean, or random attack.

The theorem assumes deterministic smooth scoring on a fixed local control-flow
branch. Discrete token or tool decisions can violate the Taylor model. Pairwise
option margins also make the primary method an offline labeled audit; open-ended
generation and code require a separately validated continuous decision score.
When inference is stochastic, clean-clean variation must be estimated and
paired with each intervention.

The threat model covers one-shot, single-edge additive interventions over a
finite unrolled horizon. It does not cover joint, persistent, adaptive
multi-round corruption or infinite-horizon stability. Relative Frobenius norms
are invariant to scalar rescaling but not arbitrary latent reparameterization,
so cross-architecture profiles require an explicit interface metric and should
not be read as intrinsic constants. Native architecture comparisons remain
confounded by different role models, training, compute, and numbers of exposed
links; matched variants are required for causal topology claims. Finally, full
access to the original question may let receivers bypass the relay. We treat
that as a measured property rather than subtracting it away.

The reported empirical results have additional, concrete limitations. The
causal and probe audits contain only 13 forced-margin clean-correct validation
examples, and the attack analysis contains only six independent examples and
18 example--handoff units. Three probe seeds quantify estimator variation but
do not triple the independent sample size. Crossed-threshold Spearman uses only
nine crossed units per seed; realized vulnerable-site top-1 uses only three
informative examples per seed, and realized site-pair concordance has only
three comparable pairs per seed. These point estimates therefore have high
sampling uncertainty and are reported without inferential confidence claims.

The PGD dose response is non-monotone because each budget is optimized
independently with one 20-step run. First-crossing analysis correctly preserves
an adversarial point found inside a nested budget, but the raw per-budget curve
is an optimizer lower bound rather than the exact robust-risk curve. The
smallest random budget also lies below the reliable direction-fidelity scale in
BF16. Finally, zero of six exact-autograd comparisons met the finite-difference
agreement rule. These facts motivate more attack restarts, a monotone
carry-forward sensitivity analysis, larger $K$, and investigation of the
autograd path, but none should be silently repaired after observing held-out
test outcomes.

\section{Conclusion}

This paper reframes latent multi-agent robustness at two resolutions: which
concrete handoff is closest to a terminal decision boundary, and what fraction
of a system's clean-correct behavior is exposed at a stated budget. \method{}
combines causal relay auditing with attack-independent random probing of
end-to-end margins, then aggregates validated handoff radii into explicit
worst-site and random-site vulnerability profiles. The construction follows the
actual unrolled computation and separates clean margin, system susceptibility,
and attack alignment. The completed validation pilot provides preliminary but
mixed support for the local diagnostic. On six examples, \method{} strongly
improves crossed-threshold ordering and small-budget flip prediction, but clean
margin remains competitive or better on other summaries. This is enough to
justify the frozen held-out test, not enough to claim general prediction or a
validated system vulnerability score. Local held-out threshold prediction
therefore remains the load-bearing empirical test; cross-dataset,
cross-architecture, prompt, and protection claims stay conditional on its
outcome.

\appendix

\section{Proofs}

\subsection{Proof of Proposition~\ref{prop:affine-radius}}

Fix a competitor \(j\) and suppress \((x,e,j)\) subscripts. Reaching its
affine boundary requires \(m+\langle g,\delta\rangle\le0\). If \(g=0\), the
strictly positive clean margin cannot be changed in the affine model. If
\(g\ne0\), Cauchy--Schwarz gives
\begin{equation}
    m\le-\langle g,\delta\rangle
    \le\|g\|_2\|\delta\|_2,
\end{equation}
so every feasible perturbation has
\(\|\delta\|_2/s\ge m/(s\|g\|_2)=m/\chi\). Equality is achieved by
\(\delta^\star=-m g/\|g\|_2^2\). The incorrect-decision region of the affine
model is the union of these competitor halfspaces, so minimizing over that
union gives \(\min_{j\ne y}m_j/\chi_j\), with the stated infinite-radius
convention.

\subsection{Proof of Proposition~\ref{prop:isotropic-recovery}}

Let $g=g_{x,e,j}$, $s=s_{x,e}$, and $q=q_e$. The exact directional derivative is
\begin{equation}
    d_j(u)=s\langle g,u\rangle.
\end{equation}
For $u$ uniform on the unit sphere in $q$ dimensions,
$\E[uu^\top]=I/q$. Therefore
\begin{align}
    \E[d_j(u)^2]
    &=s^2\E[g^\top uu^\top g]\\
    &=\frac{s^2\|g\|_2^2}{q}
    =\frac{\chi_{x,e,j}^2}{q}.
\end{align}
Rotational symmetry lets us align the first coordinate with \(g\). Using
\(\E[u_1^4]=3/[q(q+2)]\),
\begin{align}
    \operatorname{Var}\!\left(qd_j(u)^2\right)
    &=q^2\chi^4
    \left(\frac{3}{q(q+2)}-\frac{1}{q^2}\right)\\
    &=\chi^4\frac{2(q-1)}{q+2}.
\end{align}
Averaging \(K\) independent directions yields
Equation~\eqref{eq:probe-variance}; the law of large numbers gives consistency.
For the finite difference, a third-order Taylor expansion of
\(\phi(t)=m(z+tsu)\) about zero gives the standard central-difference remainder
bounded by \(Bh^2/6\).

\subsection{Proof of Theorem~\ref{thm:margin-stability}}

Fix one competitor and suppress \((x,e,j)\) subscripts. For any admissible
$\delta$, Taylor's theorem gives
\begin{equation}
    m(z+\delta)
    =m(z)+\langle \nabla m(z),\delta\rangle
    +\frac{1}{2}\delta^\top
    \nabla^2m(z+t\delta)\delta
\end{equation}
for some $t\in(0,1)$. By Cauchy--Schwarz and
Assumption~\ref{ass:curvature}, and because \(\delta=P\delta\),
\begin{align}
    m(z+\delta)
    &\ge m(z)-\|P\nabla m(z)\|_2\|\delta\|_2
    -\frac{1}{2}H\|\delta\|_2^2\\
    &\ge m(z)-\eps s\|g\|_2
    -\frac{1}{2}\eps^2s^2H,
\end{align}
which is Equation~\eqref{eq:curvature-bound}. If this lower bound is positive
for every competitor, the correct option retains the largest score. Solving the
quadratic lower bound for its positive root gives
Equation~\eqref{eq:cert-radius}, with the stated linear and constant special
cases. Truncation at \(\bar\eps\) keeps the result inside the audited ball.

\subsection{Proof of Corollary~\ref{cor:system-stability}}

If \(\eps<\min_e\rho^{\mathrm{cert}}_{x,e}\), then
\(\eps<\rho^{\mathrm{cert}}_{x,e}\) for every eligible handoff. Theorem
\ref{thm:margin-stability} therefore preserves all pairwise margins under an
intervention at any one eligible edge, proving one-shot system stability. Hence
an exact failure at budget \(\eps\) can occur only outside the set certified
robust at that budget. Taking probabilities gives
\(\Vuln^\star(\eps)\le1-C(\eps)\).

\section{Mechanistic Decomposition on the Unrolled Graph}

For a unique downstream path, submultiplicativity gives
\begin{equation}
    \chi_{x,e,j}
    \le
    s_{x,e}
    \|\nabla_{\mathrm{out}}m_{x,e,j}\|_2
    \prod_{v\in\operatorname{path}(e,\mathrm{out})}
    \|J_v\|_{\mathrm{op}}.
\end{equation}
For branching computation, the end-to-end Jacobian is a sum over paths rather
than a product of scalar role gains. The direct probe estimate remains primary
because scalar norm products discard cancellations and can be extremely loose.

\section{Preliminary Run Provenance and Evidence Status}

Table~\ref{tab:run-provenance} records the execution state underlying every
numerical result in Section~\ref{sec:results}. Input task directories and
artifacts were checked against their completion records before aggregation.
The derived six-example attack analysis is explicitly diagnostic and does not
constitute a frozen test-set result.

\begin{table}[h]
\centering
\small
\setlength{\tabcolsep}{4pt}
\caption{Provenance of the current preliminary evidence. The source hash is
the code identity recorded by the completed pilot artifacts.}
\label{tab:run-provenance}
\begin{tabular}{p{0.27\linewidth}p{0.66\linewidth}}
\toprule
Field & Recorded value \\
\midrule
Workflow root & \path{linkradius-rq2-seqlight-forcedmargin-v1} \\
Source hash &
\texttt{79815ae8cb16c6cd1478400185da4de6}\newline
\texttt{29a84be655350ae6ea50e10f5496c922} \\
System & Sequential-light on GPQA; $R=2$; batch size 1; latent length 32 \\
Decision rule & Deterministic forced-option scoring and minimum pairwise
margin; free-generation validity retained only as a diagnostic \\
Screen/clean cohort & 40 validation examples screened; 13 forced-margin
clean-correct; 13/13 stable under authenticated clean recapture \\
Causal audit & All 13 eligible examples at $P{\to}C@1$, $C{\to}S@1$, and
$S{\to}P@1$ for mismatch, identity, and zero interventions \\
Probe audit & $K=8$ antithetic directions; $h\in\{10^{-3},3\times10^{-3}\}$;
seeds 101, 202, and 303; 234 eligible configurations and 1,872 accepted pairs \\
Gradient diagnostic & Six finite example--edge rows; 0/6 finite-difference
agreements; maximum relative error 1.043; archived outside the frozen probe
gate because the rows lacked resolved system-identity metadata; not used as a
certificate \\
Attack prefix & First six of 13 eligible validation examples in frozen
execution order; 36/78 full validation attack tasks represented (three sites,
two families, six budgets), yielding 216 sample rows and 324 PGD
competitor-target rows; remaining 42 tasks not used \\
Attack settings & Per-example 20-step PGD and independent random directions;
$\eps\in\{3\!\times\!10^{-4},10^{-3},3\!\times\!10^{-3},10^{-2},
3\!\times\!10^{-2},10^{-1}\}$ \\
Data-access status & Validation only; no held-out test outcome accessed; no
architecture, prompt, protection, or independent transfer-attack result \\
\bottomrule
\end{tabular}
\end{table}

\section{Full Reproducibility Checklist}

The final appendix should report:
\begin{itemize}
    \item exact model checkpoint, precision, tokenizer, decoding rule, and
    option-scoring procedure;
    \item the chronological micro-step associated with every handoff label;
    \item shape, scale \(s_{x,e}\), basis \(B_e\), dimension \(q_e\), and
    admissible perturbation subspace of every latent object;
    \item requested and realized post-cast perturbation norms;
    \item realized probe antipodality and requested--realized cosine checks;
    \item random seeds and exact antithetic direction construction;
    \item causal-audit derangement rule and moment-matching procedure;
    \item calibration, validation, and test indices;
    \item attack training objectives, steps, restarts, and stopping rules;
    \item clean-correct and intersection clean-correct denominators;
    \item all-edge, useful-edge, worst-site, fixed-site, and random-site threat
    declarations, including validation-only link selection;
    \item treatment of censored thresholds, ties, invalid outputs, and clean
    rerun disagreement;
    \item fixed discrete traces used for theorem-bearing architecture analyses;
    \item per-example results sufficient to reproduce bootstrap intervals.
\end{itemize}

\section{Expanded Result Tables}

The final appendix should expand Table~\ref{tab:site-horizon} by dataset,
horizon, exact injection round, site, attack, and budget. It should also include
per-competitor margins, estimator seed stability, probe-radius stability,
clean-clean floors, actual perturbation norms, and all baseline comparisons.
Do not pool incompatible clean denominators or average effects across budgets
before showing the dose--response curves. System-level appendices should include
full \(\Vuln(\eps)\) curves, \(\eps_q\) values, \(\AUV\), predicted robust
accuracy, edge counts, useful-link gates, and paired architecture/prompt
comparisons on both intersection and system-specific clean sets.

\begin{thebibliography}{20}

\bibitem{du2025interlat}
Zhuoyun Du, Runze Wang, Huiyu Bai, Zouying Cao, Xiaoyong Zhu, Yu Cheng,
Bo Zheng, Wei Chen, and Haochao Ying.
\newblock \emph{Enabling Agents to Communicate Entirely in Latent Space}.
\newblock arXiv:2511.09149, 2025.
\newblock \url{https://arxiv.org/abs/2511.09149}

\bibitem{zou2025latentmas}
Jiaru Zou, Xiyuan Yang, Ruizhong Qiu, Gaotang Li, Katherine Tieu, Pan Lu,
Ke Shen, Hanghang Tong, Yejin Choi, Jingrui He, James Zou, Mengdi Wang,
and Ling Yang.
\newblock \emph{Latent Collaboration in Multi-Agent Systems}.
\newblock arXiv:2511.20639, 2025.
\newblock \url{https://arxiv.org/abs/2511.20639}

\bibitem{yang2026recursivemas}
Xiyuan Yang, Jiaru Zou, Rui Pan, Ruizhong Qiu, Pan Lu, Shizhe Diao,
Jindong Jiang, Hanghang Tong, Tong Zhang, Markus J. Buehler, Jingrui He,
and James Zou.
\newblock \emph{Recursive Multi-Agent Systems}.
\newblock arXiv:2604.25917, 2026.
\newblock \url{https://arxiv.org/abs/2604.25917}

\bibitem{wang2026outofsight}
Chenxi Wang, Ruiyang Huang, Jiayan Sun, Lei Wei, and Yifan Wu.
\newblock \emph{Out of Sight, Not Out of Mind: Unveiling Latent Attack in
Latent-based Multi-Agent Systems}.
\newblock arXiv:2605.28214, 2026.
\newblock \url{https://arxiv.org/abs/2605.28214}

\bibitem{brito2026latentagents}
Lu\'is Brito and Carlos Baquero.
\newblock \emph{When Latent Agents Lie: KV-Cache Integrity in Multi-Agent LLM
Collaboration}.
\newblock arXiv:2606.28958, 2026.
\newblock \url{https://arxiv.org/abs/2606.28958}

\bibitem{asif2026lcguard}
Sadia Asif, Mohammad Mohammadi Amiri, Momin Abbas, Prasanna Sattigeri,
and Karthikeyan Natesan Ramamurthy.
\newblock \emph{LCGuard: Latent Communication Guard for Safe KV Sharing in
Multi-Agent Systems}.
\newblock arXiv:2605.22786, 2026.
\newblock \url{https://arxiv.org/abs/2605.22786}

\bibitem{cheng2026when}
Jiaming Cheng, Subhransu Das, and Rajiv Ramnath.
\newblock \emph{When Does Latent Communication Pay? A Causal Audit of Relayed
KV Caches in Multi-Agent LLMs}.
\newblock arXiv:2608.04893, 2026.
\newblock \url{https://arxiv.org/abs/2608.04893}

\bibitem{zhang2026channels}
Huixiang Zhang and Mahzabeen Emu.
\newblock \emph{Do Latent Channels Actually Communicate? A Causal Audit of
Latent Multi-Agent LLM}.
\newblock arXiv:2607.26773, 2026.
\newblock \url{https://arxiv.org/abs/2607.26773}

\bibitem{nesterov2017random}
Yurii Nesterov and Vladimir Spokoiny.
\newblock Random gradient-free minimization of convex functions.
\newblock \emph{Foundations of Computational Mathematics}, 17(2):527--566,
2017.

\bibitem{weng2018clever}
Tsui-Wei Weng, Huan Zhang, Pin-Yu Chen, Jinfeng Yi, Dong Su, Yupeng Gao,
Cho-Jui Hsieh, and Luca Daniel.
\newblock Evaluating the robustness of neural networks: An extreme value
theory approach.
\newblock In \emph{International Conference on Learning Representations},
2018.
\newblock \url{https://arxiv.org/abs/1801.10578}

\end{thebibliography}

\end{document}
